\section{问题重述}

\subsection{问题背景}

超大规模集成电路（Very Large Scale Integration，VLSI）通过在单一芯片上集成大量电路单元，实现复杂电子系统的高度集成。随着芯片规模和设计复杂度不断提高，依靠人工方式已难以完成高效设计，电子设计自动化（Electronic Design Automation，EDA）技术因此成为现代集成电路设计的重要支撑。布图规划（Floorplanning）作为 EDA 设计流程中的核心环节之一，需要在满足模块互不重叠、不超出芯片轮廓等几何约束的条件下，合理确定各功能模块的位置与方向，为后续布局与布线奠定基础。合理的布图规划不仅需要控制芯片轮廓面积和死区比例，还需兼顾模块间互连线长，从而提高芯片面积利用率并降低后续布线开销。因此，针对不同轮廓条件和模块形状建立合理的数学模型与高效求解方法，对于提高 VLSI 物理设计质量具有重要意义。

\subsection{问题内容}

\noindent\textbf{问题一：}假设芯片轮廓不固定，且模块之间不存在连接关系。需要在允许模块旋转且互不重叠的条件下确定各模块的位置与方向，使芯片最小外接轮廓面积尽可能小；若不同方案面积相同，则优先选择长宽比更接近 $1$ 的布局。基于三组 \texttt{.blocks} 数据分别完成 \texttt{n100}、\texttt{n200} 和 \texttt{n300} 的独立布图，并给出轮廓面积、长宽比及布局可视化结果。

\noindent\textbf{问题二：}在实际固定轮廓场景下，芯片轮廓为由模块总面积和给定死区比例确定的正方形，模块之间存在网络连接关系。要求在死区比例为 $0.15$ 时，合理确定模块的位置与方向，使所有模块互不重叠且不越出芯片轮廓，并使模块间总半周长线长（HPWL）尽可能小。基于附件中的三组芯片数据，分别给出总连线长度和布局可视化结果。

\noindent\textbf{问题三：}在问题二的基础上，进一步压缩芯片轮廓，寻找仍能满足模块互不重叠且不越界条件的最小死区比例。需要分别确定三组芯片能够形成可行布局的最小死区比例，并在相应轮廓下沿用问题二的线长优化模型，重新计算总 HPWL，给出更新后的模块布局及可视化结果。

\noindent\textbf{问题四：}当功能模块由规则矩形推广为 L 型、T 型等非矩形形状，并允许进行 $90^\circ$、$180^\circ$ 和 $270^\circ$ 旋转时，需要讨论问题一中模型及求解算法的相应修正。针对题目给出的四个模块，确定其位置与旋转方向，使包围全部模块的芯片轮廓面积尽可能小，并给出最优摆放示意图及对应的最小面积。
